\section{When Can a Collision Carry an Erasure?}
\label{sec:theory}
\subsection{The exact deployed-collision set}
A quantization-conditioned attack~\cite{egashira2024} ships weights that behave well
in full precision but whose deployed form is a chosen malicious model. For erasure,
the malicious deployed model is simply the quantized original. Let $\mathcal{Q}(W)$
denote the dequantized weights the deployer obtains after recomputing every scale
from the received $W$. We call $W_*$ a \emph{scale-preserving exact collision} if it
induces the same migration and weight scales as $W_0$ and the same integer codes, so
that $\mathcal{Q}(W_*)=\mathcal{Q}(W_0)$. Such a checkpoint deploys bit-for-bit as the
quantized original; recovery is then guaranteed, and the only remaining question is
whether it can look erased in FP16.

\begin{proposition}[Pinned deployed-collision box]
\label{prop:box}
Fix the deployer's activation maxima $a$ and exponent $\alpha$. For
W8A8 with migration~\eqref{eq:smooth} and per-row INT8 weights, $W$ is a
scale-preserving exact collision with $W_0$ if and only if
(i) every column maximum $\max_r|W_{rc}|$ equals that of $W_0$;
(ii) every migrated row maximum $\max_c|W_{rc}s_c|$ equals that of $W_0$; and
(iii) every coordinate lies in its original code interval
$[(k_{rc}-\tfrac12)\sigma_r/s_c,\,(k_{rc}+\tfrac12)\sigma_r/s_c]$.
Let $\pi_0$ fix, for every row and column, one coordinate at which $W_0$ attains the
maximum. Up to rounding-tie conventions, the collisions that attain their maxima at
$\pi_0$ form the box $\mathcal{B}_{\pi_0}$: the code intervals intersected with the
column caps $|W_{rc}|\le\max_{r'}|W_{0,r'c}|$ and row caps $|W_{rc}|\le127\sigma_r/s_c$,
with the coordinates of $\pi_0$ held at their original values. For weight-only
round-to-nearest quantizers with per-row or per-group absmax scales (INT8, INT4,
NF4~\cite{qlora}), the same holds with group caps and group arg-max coordinates in
place of (i)--(ii).
\end{proposition}
\begin{proof}
By~\eqref{eq:smooth}, $s_c$ depends on $W$ only through the column maximum, and
$\sigma_r$ only through the migrated row maximum. Conditions (i)--(ii) are therefore
equivalent to equality of all scales, and given equal scales, (iii) is equality of
integer codes. The caps prevent a maximum from growing, and holding the coordinates of
$\pi_0$ fixed keeps each maximum attained, so every point of $\mathcal{B}_{\pi_0}$ is a
collision; conversely a collision attaining its maxima at $\pi_0$ satisfies every
constraint of $\mathcal{B}_{\pi_0}$. Weight-only quantizers have a single absmax scale
per group, giving the analogous box.
\end{proof}

The full collision set is the union of $\mathcal{B}_\pi$ over all admissible arg-max
patterns $\pi$, and it need not be convex. For example, with $a=(1,1)$ and
$\alpha=0.5$, the matrices with rows $(1,0.999),(0.999,1)$ and $(0.999,1),(1,0.999)$ are
both collisions, but their midpoint has smaller column maxima and is not. All
constructions and certificates below use the original pattern $\pi_0$, so their
guarantees are statements about $\mathcal{B}_{\pi_0}$, not about the whole union.

Condition (i) is easy to miss. Our earlier construction capped rows but not columns,
so recalibrating on the crafted checkpoint moved the migration scales and code
agreement plateaued at 0.83--0.84 after 16 iterations (Appendix~\ref{app:earlier}).
With column caps, agreement after the deployer recalibrates is 1.0 for every tensor
(Section~\ref{sec:feasibility}). Because the box is defined by the deployer's own
statistics, a bound over $\mathcal{B}_{\pi_0}$ also covers attackers who do not know
the calibration data but still produce a collision with pattern $\pi_0$. Attacks that instead \emph{change} the scales,
such as outlier injection~\cite{zhan2026}, deploy to a different model and are
outside this class (Section~\ref{sec:limits}).

\subsection{Key/value outputs are a sufficient statistic}
UCE and TIME modify only the cross-attention key and value maps
$W\in\mathbb{R}^{d_o\times768}$. Their inputs are the text encoder's 77 token
embeddings for the prompt and, under classifier-free guidance, for the empty prompt.
Neither depends on the UNet, the latent, or the timestep. For a fixed prompt, seed,
and scheduler, a generated image therefore depends on $W$ \emph{only through} the
products $We$ over these embeddings. Two checkpoints with equal key/value outputs on
a prompt produce identical images.

The behavioral effect of an edit $\Delta=W_E-W_0$ is thus fully determined by its
output displacement $\Delta E$, where $E$ stacks the relevant embeddings. We measure
how closely a collision in $\mathcal{B}_{\pi_0}$ can reproduce this displacement:
\begin{equation}
 \rho^*=\min_{X\in\mathcal{B}_{\pi_0}-W_0}\frac{\|(X-\Delta)E\|_F}{\|\Delta E\|_F}.
 \label{eq:capacity}
\end{equation}
This is a convex, row-separable quadratic program in the Gram metric $G=EE^\top$.
Projecting $\Delta$ coordinate-wise onto the box, as in our earlier construction and
in the repair step of prior quantization attacks, corresponds to $G=I$ and ignores
that text embeddings are highly anisotropic. We solve~\eqref{eq:capacity} for each
tensor with FISTA under a diagonal majorizer, and certify it with the Frank--Wolfe
duality gap. The gap yields a \emph{lower} bound on $\rho^*$: no point of
$\mathcal{B}_{\pi_0}$ reproduces this edit's key/value displacement on the fit prompts
more closely, whatever optimizer is used. We call the minimizer the \emph{optimal
reconstruction attack}. The bound concerns reconstruction of this particular edit, not
behavior: a checkpoint could in principle suppress the concept through a different
key/value displacement.

\paragraph{From outputs to behavior}
We report the explained fraction $f=\langle XE,\Delta E\rangle/\|\Delta E\|_F^2$ on
held-out prompts. Full reproduction ($f\approx1$, small residual) implies FP16
behavior close to the erasure by sufficiency. Partial reproduction carries no such
implication. The honest ladder $W_0+t\Delta$ shows that suppression is a nonlinear
function of the edit, and an attack's orthogonal residual can also matter. We
therefore use~\eqref{eq:capacity} to construct the attack and to explain its outcome,
and judge stealth only from measured behavior (Section~\ref{sec:behavior}).
